\emph{संपूर्णता प्रमेय} हे \emph{निर्दोषता प्रमेयाचे} उलट विधान आहे. एका मांडणीनुसार $\Gamma \Entails !A$ असेल, तर $\Gamma \Proves !A$. दुसऱ्या मांडणीनुसार $\Gamma$ सुसंगत असेल, तर तो पूर्ततायोग्य असतो. आपण दुसरी मांडणी सिद्ध केली (आणि तिच्यावरून पहिली मिळवली). ही सिद्धता गुंतागुंतीची आहे आणि तिच्यात अनेक टप्पे आहेत. आपण सुसंगत वाक्यसंच~$\Gamma$ पासून सुरुवात करतो. आधी भाषेत अनंत नवीन स्थिरचिन्हे~$c_i$ घालतो. विस्तारित भाषेतील फक्त~$x$ हा चर मुक्त असलेल्या प्रत्येक सूत्र~$!A(x)$ साठी $\lexists[x][!A(x)] \lif !A(c)$ या स्वरूपाचे वाक्य जोडतो. त्यासाठी असे नवीन स्थिरचिन्ह~$c$ निवडतो की ते~$!A(x)$ मध्येही येत नाही आणि तोपर्यंतच्या वाक्यसंचातही येत नाही. यातून~$\Gamma$ चा समावेश असलेला \emph{संतृप्त} सुसंगत वाक्यसंच मिळतो. आता त्या संचाचा विस्तार करून \emph{संपूर्ण सुसंगत संच} मिळवतो. संपूर्ण सुसंगत संचाचा असा गुणधर्म आहे की कोणतेही वाक्य~$!A$ घेतले, तर $!A$ किंवा $\lnot !A$ यांपैकी एक त्या संचात असते (दोन्ही कधीही नसतात). सुरुवातीचा संच संतृप्त असल्याने आता आपल्याकडे~$\Gamma$ चा समावेश असलेला संतृप्त, संपूर्ण, सुसंगत वाक्यसंच~$\Gamma^*$ आहे. या संचावरून अशी रचना~$\Struct{M(\Gamma^*)}$ परिभाषित करता येते की प्रत्येक वाक्य~$!A$ साठी $\Sat{M(\Gamma^*)}{!A}$ तेव्हा आणि केवळ तेव्हाच, जेव्हा $!A \in \Gamma^*$. विशेषतः $\Sat{M(\Gamma^*)}{\Gamma}$, म्हणजे $\Gamma$ पूर्ततायोग्य आहे. $=$ हे चिन्ह असेल, तर रचना घडवण्याची ही पद्धत थोडी अधिक गुंतागुंतीची असते.

संपूर्णता प्रमेयावरून दोन महत्त्वाची अनुप्रमेये मिळतात. \emph{संहतता प्रमेयानुसार}, $\Gamma \Entails !A$ तेव्हा आणि केवळ तेव्हाच, जेव्हा असा सांत उपसंच~$\Gamma_0 \subseteq \Gamma$ असतो की $\Gamma_0 \Entails !A$. याची समतुल्य मांडणी अशी आहे की प्रत्येक सांत $\Gamma_0 \subseteq \Gamma$ पूर्ततायोग्य असेल तेव्हा आणि केवळ तेव्हाच $\Gamma$ पूर्ततायोग्य असतो. संहतता प्रमेय विशिष्ट गुणधर्म असलेल्या रचनांचे अस्तित्व सिद्ध करण्यासाठी उपयोगी पडते. उदाहरणार्थ, ज्या उपपत्तीला हवी तितकी मोठी सांत प्रतिमाने असतात तिला अनंत प्रतिमानेही असतात, हे त्यावरून दाखवता येते. विशेषतः, याचा अर्थ असा की सांतता प्रथम-क्रम तर्कशास्त्रात व्यक्त करता येत नाही. दुसरे अनुप्रमेय, \emph{लोव्हेनहाइम--स्कोलेम प्रमेय}, असे सांगते की गणनीय भाषेतील वाक्यांचा संच $\Gamma$ पूर्ततायोग्य असेल, तर त्याला गणनीय प्रतिमान असते. त्यावरून अगणनीयता प्रथम-क्रम तर्कशास्त्रात व्यक्त करता येत नाही हे दिसते.
